% Fragmento público generado desde propietarios íntegros.
% Residencia: 52.6.3.
% Fuente: ../incoming/radion_monografia_20260827/manuscrito/sections/07_incertidumbre_oscilador.tex:112-155
\subsubsection{Oscillateur inductif--capacitif associé au couple angulaire}

\paragraph{Construction indépendante des 2 réservoirs conjugués}

Étant donnés \(\omega>0\) et une impédance de référence \(Z_*>0\), nous définissons
\begin{equation}
L_\eta=\frac{Z_*}{\omega}e^{-\eta},
\qquad
C_\eta=\frac{1}{Z_*\omega}e^\eta.
\label{eq:LC-def}
\end{equation}
Alors
\begin{equation}
L_\eta C_\eta=\omega^{-2},
\qquad
\sqrt{\frac{L_\eta}{C_\eta}}=Z_*e^{-\eta}.
\label{eq:LC-invariants}
\end{equation}
La fréquence reste fixe ; l'impédance conserve l'orientation de la
forme.

L'hamiltonien est
\begin{equation}
H_{LC}=\frac{q^2}{2C_\eta}+\frac{\Phi^2}{2L_\eta}.
\label{eq:HLC}
\end{equation}
Avec les variables normalisées
\begin{equation}
Q=\sqrt{Z_*}\,q,\qquad
P=\frac{\Phi}{\sqrt{Z_*}},
\label{eq:LC-map}
\end{equation}
on obtient
\begin{equation}
H_{LC}=\frac\omega2
\left(e^{-\eta}Q^2+e^\eta P^2\right).
\label{eq:HLC-QP}
\end{equation}
Sur \eqref{eq:elipse-param},
\begin{equation}
H_{LC}=\hret\omega.
\label{eq:LC-level}
\end{equation}

% Fuente: ../incoming/radion_monografia_20260827/manuscrito/sections/07_incertidumbre_oscilador.tex:198-239
\paragraph{Oscillateur LC comme unité dynamique élémentaire}

L'oscillateur harmonique traverse presque toute la physique parce qu'il linéarise le
voisinage d'un équilibre, décompose les champs en modes, organise les quanta et
permet de lire la propagation comme une superposition de phases. Dans le radion, cette
universalité reçoit une généalogie concrète : le mode ne commence pas par une
équation différentielle externe, mais par une forme positive dont l'action de tour est
fixe. La dynamique conventionnelle apparaît lorsqu'on choisit une horloge \(\omega\) et un
transducteur \(Z_*\).

La fréquence ne modifie pas l'ellipse d'action : elle change l'énergie du niveau
au moyen de \(E=\hret\omega\). Cette séparation est cruciale. La géométrie fixe
l'action ; l'horloge convertit l'action en énergie.

\subsubsection{Réversion temporelle et propagations conjuguées}

Soit \(H=H^*\) et soit \(J_E\) une structure complexe avec \(J_E^2=-I\).
Supposons un opérateur involutif \(C\) tel que
\[
CJ_EC^{-1}=-J_E,\qquad CHC^{-1}=H.
\]
L'évolution
\begin{equation}
U(t)=\exp\left(-\frac{J_EHt}{\hret}\right)
\label{eq:evolution-J}
\end{equation}
satisfait exactement
\begin{equation}
CU(t)C^{-1}=U(-t).
\label{eq:time-reversal}
\end{equation}
Telle est la forme opératoire minimale de la bidirectionnalité : la conjugaison
échange les prolongations temporellement opposées.

La théorie de l'absorbeur de Wheeler--Feynman et la lecture des antiparticules
comme trajectoires temporellement conjuguées offrent une réalisation historique
ultérieure. L'identité \eqref{eq:time-reversal} est démontrée ; la transformer
en une théorie complète des propagateurs avancés et retardés exige de déclarer
l'action, les conditions de frontière, la causalité et les interactions de l'absorbeur.
La monographie conserve l'intuition physique sans transformer une correspondance
en une égalité de théories.

